\section{Soundness Theorem}\label{sec:soundness}

\subsection{Statement}

\begin{theorem}[Soundness of BEAMFS recovery]\label{thm:soundness}
Let~\(\Lambda\) be a layout satisfying invariants~\textup{I1--I3} of
\cref{sub:invariants}. Let~\(\mathcal{G}\) be the recovery operator
of \cref{def:G}, satisfying~\textup{I4--I5}. Let~\(\mathcal{E}\) be a
family of physical perturbations such that, for every event
\(e \in \mathcal{E}\) and every state~\(s\), the number of corrupted
symbols in any single RS subblock of~\(e(s)\) does not exceed the
correction capacity~\(t = \lfloor (n-k)/2 \rfloor\) declared by~\(\Lambda\).
Then for every initial consistent state~\(s_0 \models\) and every
\(e \in \mathcal{E}\), there exists~\(K \le |\Lambda|\) such that
\[
  \mathcal{G}^{K}\bigl(e(s_0)\bigr) \models
  \quad\text{or}\quad
  \mathcal{G}^{K}\bigl(e(s_0)\bigr) = \bot,
\]
where~\(|\Lambda|\) denotes the number of regions in~\(\Lambda\).
\end{theorem}

The full proof appears in \cref{app:proof}. The argument has
three steps:

\paragraph{Termination.}
By Lemma~\ref{lem:descent}, the function~\(N\) strictly decreases on
each non-fail-closed iteration. Since~\(N \in \{0, 1, \dots, |\Lambda|\}\)
and~\(N = 0\) implies consistency by \cref{def:consistent},
the recursion terminates in at most~\(|\Lambda|\) iterations.

\paragraph{Safety.}
At every iteration, if step~3 or step~5 of \cref{def:G}
fails (CRC mismatch after RS, or sentinel violation),~\(\bot\) is
returned. By \cref{prop:absorb}, fail-closed is absorbing:
no subsequent iteration can return a state that silently masks the
violation.

\paragraph{Coverage of the perturbation family.}
The hypothesis on~\(\mathcal{E}\) (no event corrupts more than~\(t\)
symbols in any single subblock) is exactly the RS code's correction
capacity. Events that violate this hypothesis correspond, by~I1, to
physical perturbations whose intensity exceeds the layout's design
target. They are out of scope of the theorem and of the layout, and
the operator's response (fail-closed at step~3 or~5) is the
designed safe failure.

\subsection{Discussion}

\paragraph{What this theorem does, and does not, prove.}
\Cref{thm:soundness} establishes that the recovery operator either
returns to consistency or to fail-closed in a bounded number of
steps, under stated physical hypotheses. It does \emph{not} prove
that the kernel implementation of~\(\mathcal{G}\) is bug-free: a
buggy implementation can violate I4 or I5 even with a correct
operator definition. \Cref{sec:m2c-tier3} sketches an implementation
path that preserves the invariants by construction, but only a
verified compilation chain (in the seL4~\cite{klein2009sel4} or
FSCQ~\cite{chen2015fscq} sense) can close the implementation gap.
We mark this as v3 work.

\paragraph{Why the bound is~\(|\Lambda|\), not a constant.}
The bound scales with the number of regions because the worst case is
a perturbation that corrupts every region simultaneously up to the
RS limit. This worst case is unrealistic (\cref{sec:phys} establishes
that simultaneous corruption of all regions requires fluences
orders of magnitude beyond environmental backgrounds), but the bound
is conservative and has the merit of making no probabilistic
assumption. The per-iteration cost is dominated by a single
RS-decode of~\(O(n^{2})\) GF operations
(\cref{sec:m2c-tier2}), so the total worst-case running time is
\(O(|\Lambda| \cdot n^{2})\): sub-millisecond on contemporary CPUs
for the FTRFS layout.

\paragraph{Comparison with crash-consistency proofs.}
FSCQ~\cite{chen2015fscq} proves that a filesystem returns to a
consistent state after a crash (sudden power loss). The hypothesis
there is that the storage substrate is \emph{benevolent}: it
preserves whatever was written, exactly. \Cref{thm:soundness}
proves the dual: it assumes the substrate is \emph{adversarial}
within the bounded-perturbation family~\(\mathcal{E}\), and the
filesystem must recover. The two proofs compose: a BEAMFS-compliant
implementation satisfying \cref{thm:soundness} \emph{and} satisfying
crash-Hoare logic in the FSCQ sense recovers from the union of
both threat models. Combining the two formal frameworks is left
for future work.
